\documentclass[a4paper,
fontsize=11pt,
%headings=small,
oneside,
numbers=noperiodatend,
parskip=half-,
bibliography=totoc,
final
]{scrartcl}

\usepackage[babel]{csquotes}
\usepackage{synttree}
\usepackage{graphicx}
\setkeys{Gin}{width=.4\textwidth} %default pics size

\graphicspath{{./plots/}}
\usepackage[ngerman]{babel}
\usepackage[T1]{fontenc}
%\usepackage{amsmath}
\usepackage[utf8x]{inputenc}
\usepackage [hyphens]{url}
\usepackage{booktabs} 
\usepackage[left=2.4cm,right=2.4cm,top=2.3cm,bottom=2cm,includeheadfoot]{geometry}
\usepackage[labelformat=empty]{caption} % option 'labelformat=empty]' to surpress adding "Abbildung 1:" or "Figure 1" before each caption / use parameter '\captionsetup{labelformat=empty}' instead to change this for just one caption
\usepackage{eurosym}
\usepackage{multirow}
\usepackage[ngerman]{varioref}
\setcapindent{1em}
\renewcommand{\labelitemi}{--}
\usepackage{paralist}
\usepackage{pdfpages}
\usepackage{lscape}
\usepackage{float}
\usepackage{acronym}
\usepackage{eurosym}
\usepackage{longtable,lscape}
\usepackage{mathpazo}
\usepackage[normalem]{ulem} %emphasize weiterhin kursiv
\usepackage[flushmargin,ragged]{footmisc} % left align footnote
\usepackage{ccicons} 
\setcapindent{0pt} % no indentation in captions
\usepackage{xurl} % Breaks URLs

%%%% fancy LIBREAS URL color 
\usepackage{xcolor}
\definecolor{libreas}{RGB}{112,0,0}

\usepackage{listings}

\urlstyle{same}  % don't use monospace font for urls

\usepackage[fleqn]{amsmath}

%adjust fontsize for part

\usepackage{sectsty}
\partfont{\large}

%Das BibTeX-Zeichen mit \BibTeX setzen:
\def\symbol#1{\char #1\relax}
\def\bsl{{\tt\symbol{'134}}}
\def\BibTeX{{\rm B\kern-.05em{\sc i\kern-.025em b}\kern-.08em
    T\kern-.1667em\lower.7ex\hbox{E}\kern-.125emX}}

\usepackage{fancyhdr}
\fancyhf{}
\pagestyle{fancyplain}
\fancyhead[R]{\thepage}

% make sure bookmarks are created eventough sections are not numbered!
% uncommend if sections are numbered (bookmarks created by default)
\makeatletter
\renewcommand\@seccntformat[1]{}
\makeatother

% typo setup
\clubpenalty = 10000
\widowpenalty = 10000
\displaywidowpenalty = 10000

\usepackage[colorlinks, linkcolor=black,citecolor=black, urlcolor=libreas,
breaklinks= true,bookmarks=true,bookmarksopen=true]{hyperref}
\usepackage{breakurl}
\usepackage{hyperxmp}
\usepackage{accessibility}

%meta
%meta

\fancyhead[L]{E. Bunge\\ %author
LIBREAS. Library Ideas, 48 (2026). % journal, issue, volume.
\href{https://doi.org/10.18452/37906}{\color{black}https://doi.org/10.18452/37906}
{}} % doi  
\fancyhead[R]{\thepage} %page number
\fancyfoot[L] {\ccLogo \ccAttribution\ \href{https://creativecommons.org/licenses/by/4.0/}{\color{black}Creative Commons BY 4.0}}  %licence
\fancyfoot[R] {ISSN: 1860-7950}

\title{\LARGE{Das große Krabbeln in der Bibliothek: ein Einblick in die Integrierte Schädlingsbekämpfung in Kultureinrichtungen und den Einfluss des Klimawandels}}% title
\author{Eva Bunge} % author

\setcounter{page}{1}

\hypersetup{%
      pdftitle={Das große Krabbeln in der Bibliothek: ein Einblick in die Integrierte Schädlingsbekämpfung in Kultureinrichtungen und den Einfluss des Klimawandels},
      pdfauthor={Eva Bunge},
      pdfcopyright={CC BY 4.0 International},
      pdfsubject={LIBREAS. Library Ideas, 48 (2026)},
      pdfkeywords={Integrierte Schädlingsbekämpfung, Klimawandel, Kulturgüterschutz, integrated pest management, climate change, cultural heritage protection},
      pdflicenseurl={https://creativecommons.org/licenses/by/4.0/},
      pdfcontacturl={http://libreas.eu},
      pdfurl={https://doi.org/10.18452/37906},
      pdfdoi={10.18452/37906},
      pdflang={de},
      pdfmetalang={de}
     }



\date{}
\begin{document}

\maketitle
\thispagestyle{fancyplain} 

%abstracts
\begin{abstract}
\noindent
\textbf{Kurzfassung}: Wir alle mussten uns wohl schon einmal mit ihnen
auseinandersetzen, auch wenn die wenigsten von uns gerne darüber
sprechen: Mehlwürmer, Kleidermotten, Ameisen, Mäuse, Silber- und
Papierfischchen, Schaben, Teppichkäfer und ähnliches Getier, das in der
Wohnung umherkreucht und fleucht. Zusammengefasst werden diese Tiere oft
als Schädlinge bezeichnet. Denn auch wenn sie biologisch betrachtet
durchaus wichtige Funktionen wahrnehmen, so richten sie in und an \linebreak
menschlichen Behausungen doch oft beträchtlichen Schaden an.
Kultureinrichtungen wie Archive, Bibliotheken und Museen sind von dieser
Gefahr nicht ausgenommen und haben aufgrund des in ihnen bewahrten, oft
einzigartigen Kulturguts eine besondere Sorgfaltspflicht auch
hinsichtlich der Schädlingsbekämpfung. Dieser Beitrag soll anhand einer
kleinen Wanderausstellung, die vom Januar bis April 2026 am Deutschen
Museum in München Station machte, einen Überblick über Hintergründe,
Methoden und aktuelle Entwicklungen der sogenannten \enquote{Integrierten
Schädlingsbekämpfung} (oft wird auch im Deutschen der englische
Fachbegriff IPM für \enquote{Integrated Pest Management} verwendet) geben.

\begin{center}\rule{0.5\linewidth}{0.5pt}\end{center}

\noindent\textbf{Abstract}: Most of us have likely had to deal with them at some
point, even if few of us enjoy talking about it: mealworms, clothes
moths, ants, mice, silverfish and paperfish, cockroaches, carpet
beetles, and similar critters crawling and scurrying around the home.
Collectively, these creatures are often referred to as pests. Although
they perform vital biological functions in nature, they often cause
considerable damage within and to human dwellings. Cultural institutions
such as archives, libraries, and museums are not exempt from this threat
and, due to the often unique cultural heritage they preserve, bear a
special responsibility regarding pest control. Based on a small
traveling exhibition that was hosted at the Deutsches Museum in Munich
from January to April 2026, this article provides an overview of the
background, methods, and current developments in the so called
\enquote{Integrierte Schädlingsbekämpfung} (often referred to by the English
term IPM for \enquote{Integrated Pest Management}, even in German).
\end{abstract}

%body
\newpage
Als die kleine Kabinettausstellung \enquote{Hier nagt nicht nur der Zahn
der Zeit -- Museumskäfer, Motten und Schimmel in Zeiten des
Klimawandels}, erstellt vom Naturhistorischen Museum Wien, im Januar
2026 im Verkehrszentrum des Deutschen Museums in München eröffnete,
wurde in der begleitenden Pressemitteilung die Leiterin der
Restaurierungswerkstätten des Münchner Museums zitiert mit: \enquote{Da
haben wir den Kofferraum des Autos aufgemacht, und es kam uns eine ganze
Wolke von Motten entgegen.}\footnote{Deutsches Museum (13.01.2026).
  \enquote{Die Motten sind unser größter Feind},
  \url{https://www.deutsches-museum.de/museum/aktuell/die-motten-sind-unser-groesster-feind}
  {[}zuletzt aufgerufen 12.04.2026{]}.} Dass man so offen über die
eigenen Probleme mit Schädlingsbefall spricht und damit auch preisgibt,
dass das anvertraute Kulturgut Schaden genommen hat, ist nicht
selbstverständlich und im ersten Moment vielleicht auch etwas
unangenehm. Zumindest im Museumsbereich wird es aber immer
selbstverständlicher, auch darüber zu kommunizieren. So kann die
Öffentlichkeit über die Arbeit der Kulturinstitutionen, mit all den
zugehörigen Schwierigkeiten und Aufwänden, aufgeklärt und auch ein Bezug
zu den Belangen von Privatpersonen hergestellt werden.
Schädlingsbekämpfung, die Insekten, an die man in diesem Kontext meist
zuerst denkt, und das laufende Insektensterben, sind alltägliche
Berührungspunkte, die durch zunehmende Globalisierung und Klimawandel
immer weiter in unser Leben eingreifen und an Bedeutung gewinnen.

Die angesprochene Ausstellung leistete einen kleinen Beitrag zu dieser
Kommunikation. Sie gliederte sich in drei Teile, an denen sich auch der
Aufbau dieses Beitrags orientiert: Zuerst ein Einblick in häufige
Museums- und Wohnungsschädlinge, daran anschließend historische und
aktuelle Bekämpfungsmethoden im Überblick und abschließend wird der
Einfluss des Klimawandels auf Schädlinge thematisiert.


\begin{figure}
\centering
\captionsetup{width=0.9\textwidth}
\includegraphics[width=0.9\textwidth]{img/abb1.jpg}
\caption{Abbildung 1: Einer der drei Schaukästen der Wanderausstellung gibt einen
Einblick in Schädlingsbefall und seine Folgen.}
\alt{Ein blassgelber, hölzerner Ausstellungsschrank: Oben die Überschrift "Im Küchenschrank und im Museumsdepot: Altbekannte Schädlinge". Darunter sind mehrere Ausstellungsstücke zu sehen, insbesondere ausgestopfte und zerfressene Präparate einer Eule und einer Katze, sowie befallene Insekten- und Pflanzenkästen, Stoffreste und ein Buch.}
\end{figure}

\section{Museums- und
Wohnungsschädlinge}\label{museums--und-wohnungsschuxe4dlinge}

Dass ein Technikmuseum wie das Verkehrszentrum des Deutschen Museums --
in dem vor allem Autos, Flugzeuge, Eisenbahnen und ähnliche Gefährte
ausgestellt sind -- stark mit Schädlingen zu kämpfen hat, mag auf den
ersten Blick verwundern. Bei genauerem Hinsehen erklärt sich dieses
Phänomen aber schnell: Schädlinge benötigen zur Ernährung organisches
Material. Neben Lebensmitteln zählen dazu beispielsweise auch Leder,
Fell, Holz, Stroh, Leim, Papier sowie Web- und Füllstoffe aus Wolle,
Baumwolle, Leinen und ähnlichen Rohstoffen. In älteren Fahrzeugen wurden
Kopfstützen, Sitze, Gestelle oder eben der Teppichbelag im Kofferraum
eines Autos aus eben diesen Materialien gefertigt und bieten damit
besten Boden für allerlei Insekten und Nagetiere.

Noch viel mehr bedroht sind naturhistorische und ethnologische
Sammlungen, aber auch \linebreak Kunstmuseen, in denen präparierte Tiere und
menschliches Kunsthandwerk bewahrt werden. Schädlinge sind hier schon
ein altbekanntes Problem, was man exemplarisch am sogenannten
Museumskäfer (\emph{Anthrenus museorum}), einer Unterart der Speckkäfer
(\emph{Dermestidae}), festmachen kann, der schon 1761 von Carl von Linné
diesen Namen erhielt und von ihm mit \enquote{Habitat in museis, quæ
destruit}\footnote{Carl von Linné (1761). Fauna Svecica. Stockholm: Lars
  Salvius, S. 145.
  \url{https://www.biodiversitylibrary.org/page/32170614} {[}zuletzt
  aufgerufen 12.04.2026{]}.} beschrieben wurde. Er frisst Keratin, einen
Stoff, der in Haaren und Chitin enthalten ist, und fühlt sich somit in
biologischen Sammlungen pudelwohl.

Ein neuartiger Schädlingsbefall, der letztes Jahr durch die Presse ging,
ist das Auftreten von Schimmelpilzen der Sorte \emph{Aspergillus
restrictus} in einer ganzen Reihe von dänischen Museen.\footnote{Miranda
  Bryant (06.05.2025). Denmark's museum objects at risk from
  \enquote*{extreme} new mould, say conservators. In: The Guardian,
  \url{https://www.theguardian.com/world/2025/may/06/denmarks-museum-objects-at-risk-from-extreme-new-mould-say-conservators}
  {[}zuletzt abgerufen 12.04.2026{]}.} Diese Pilze zeichnen sich vor
allem dadurch aus, dass sie im Unterschied zu den bekannteren,
Feuchtigkeit liebenden Schimmelpilzen, ein trockenes Raumklima
bevorzugen und sich entsprechend auch in klimatisierten und alle
aktuellen Anforderungen des Bestandsschutzes erfüllenden Museumsdepots
ausbreiten konnten. Hier werden in Zukunft neue Wege gefunden werden
müssen, um Museumsobjekte, aber auch Archivalien und Buchbestände zu
schützen.

Einen Überblick über alle relevanten Schädlinge zu geben, ist an dieser
Stelle sicherlich nicht möglich. Aber einige besonders für Bibliotheken
relevante Insekten sollen an dieser Stelle kurz aufgegriffen werden.
Insbesondere ist hier das Papierfischchen (\emph{Ctenolepisma
longicaudatum}) zu nennen, der größere Verwandte der bekannten
Silberfischchen. Während diese sich gern in feuchten Bädern und Küchen
an Milben, Schimmel und anderem gütlich tun, kommen Papierfischchen auch
mit einem wesentlich trockeneren und wärmeren Raumklima zurecht. Sie
können bis zu circa sechs Jahre alt werden, sind anspruchslos,
widerstandsfähig und ernähren sich, wie der Name sagt, von Papier und
Pappe. Ihre Fraßspuren tauchen als unregelmäßige, mäandernde Gänge im
Papier auf und können bei starkem Befall signifikante Schäden an
Buchbeständen verursachen. Bemerkenswert am Papierfischchen ist, dass es
vor etwa zehn bis fünfzehn Jahren kaum Thema im Fachdiskurs der
europäischen Schädlingsbekämpfung war, sich aber seitdem rasant
ausgebreitet hat -- nicht zuletzt durch den florierenden Online-Handel,
mit dessen Paketen es um die ganze Welt reist.

Der Brotkäfer (\emph{Stegobium paniceum}), ein Vertreter der Nagekäfer,
frisst nicht nur Brot, sondern auch gern den stärkehaltigen Kleber, der
in alten Bucheinbänden verwendet wurde. Obwohl er nur wenige Millimeter
klein ist, kann er in Massen auftreten und unwiderrufliche Zerstörungen
verursachen -- so geschehen beispielsweise in den Bibliotheken des
Wiener Schottenstifts oder der zum UNESCO-Weltkulturerbe gehörenden
Territorialabtei Pannonhalma in Ungarn.\footnote{Tanja Beuthien (2026).
  Motten? Alarmstufe rot! In: GEO 03/26, S. 48--64.}

Früher in Wohnungen und Kultureinrichtungen weit verbreitet und heute
aufgrund des wärmeren und trockeneren Raumklimas seltener geworden, ist
die umgangssprachlich als Holzwurm bekannte Larve des Gemeinen
Nagekäfers (\emph{Anobium punctatum}). Er ist heute noch in feuchten
Kellern und Gemäuern zu finden, wird in beheizten Räumen aber langsam
durch den aus dem Mittelmeerraum kommenden Südlichen Nagekäfer
(\emph{Oligomerus ptilinoides}) ersetzt, der sich auch in trockenerem
Holz wohlfühlt und sich in hölzernen Buchdeckeln und alten Regalen und
Fußböden ausbreiten kann.

Neben den genannten Arten existiert noch eine fast unüberschaubare und
sich stetig wandelnde Vielzahl an potentiellen Schädlingen: Speck-,
Teppich-, Pelz-, Museums-, Nage-, Wespen- und Wodkakäfer, aber auch
Silber-, Papier- und Geisterfischchen sowie Kleider-, Pelz- und
Lebensmittelmotten und viele mehr stellen eine potentielle Gefahr für
Bibliotheken und andere Kultureinrichtungen dar. Inzwischen wurden
ausgeklügelte Systeme entwickelt, um durch sie verursachte Schäden zu
minimieren.

\section{Historische und aktuelle
Bekämpfungsmethoden}\label{historische-und-aktuelle-bekuxe4mpfungsmethoden}

Um Schädlingen und insbesondere Insekten vorzubeugen oder aktive Befälle
zu bekämpfen, wurden schon in der ersten Hälfte des letzten Jahrhunderts
verschiedenste Chemikalien eingesetzt, die in den allermeisten Fällen
nicht nur für Tiere, sondern auch für Menschen äußerst
gesundheitsschädlich sind. Chloroform, Arsen und
Dichlordiphenyltrichlorethan, das unter der Abkürzung DDT Aufsehen
erregte, sind nur einige der bekannteren eingesetzten Gifte. Ihre
Wirkung hallt noch heute nach, denn viele Museumssammlungen wurden so
intensiv mit ihnen behandelt, dass einige Depots heute nur noch mit
Schutzkleidung betreten werden können. Die Gefahr geht dabei nicht nur
von den Objekten selbst aus, sondern beispielsweise auch von dem Staub,
der durch sie kontaminiert werden kann und über Hautkontakt oder die
Atemwege den Weg in den menschlichen Körper findet.\footnote{Elise
  Spiegel et al.~(2019). Handreichung zum Umgang mit kontaminiertem
  Sammlungsgut. München: oekom Verlag.}

Heute gehen Kultureinrichtungen mit dem Konzept der
\enquote{Integrierten Schädlingsbekämpfung} (IPM) einen nachhaltigeren,
umwelt- und gesundheitsfreundlicheren Weg der Bestandserhaltung. Bereits
2016 in der Norm \enquote{DIN EN 16790:2016-12 Erhaltung des kulturellen
Erbes -- Integrierte Schädlingsbekämpfung (IPM) zum Schutz des
kulturellen Erbes} festgehalten und standardisiert, baut es auf den drei
Säulen Monitoring/Diagnose, Prävention und Schädlingsbekämpfung
auf.\footnote{Maria Kobold und Jana Moczarski (2019). Bestandserhaltung:
  ein Ratgeber für Verwaltungen, Archive und Bibliotheken. 3.,
  überarbeitete und erweiterte Auflage. Darmstadt: Universitäts- und
  Landesbibliothek Darmstadt, S. 82 ff.,
  \url{https://nbn-resolving.org/urn:nbn:de:tuda-tuprints-114079}
  {[}zuletzt abgerufen 13.042026{]}.}

Um die Eckpfeiler des IPM zu illustrieren, sollen hier beispielhaft die
Maßnahmen am Deutschen Museum (DM) in München skizziert werden. Sie
decken sicherlich nicht alle Aspekte des IPM umfassend ab, geben aber
einen guten Einblick in die wichtigsten Anwendungen. So wird das IPM am
DM in allen Sammlungsbereichen, also in Archiv, Bibliothek und
Objektsammlung, angewendet. Zum Monitoring werden verschiedenste Fallen
ausgebracht, die bedarfsgerecht mehrmals pro Jahr ausgewechselt und
ausgewertet werden. Durch das Führen einer detaillierten Statistik,
ausgewertet nach Zeitpunkt, Ort, Art und Häufigkeit der gefundenen
Insekten, können Rückschlüsse auf allgemeine Trends der
Schädlingsentwicklung, aber auch akute Befälle mit dringendem
Handlungsbedarf identifiziert werden.

Je nach Art der überwachten Sammlungen kommen dabei unterschiedliche
Fallen zum Einsatz. In Archiv und Bibliothek sind dies vor allem
dreieckige Klebefallen aus Pappe, die auf dem Fußboden und in den
Regalen ausliegen. Zur Überwachung der besonders relevanten
Papierfischchen ist nicht einmal das Anbringen eines Lockstoffs
erforderlich, da sie von dem Pappkarton angezogen werden, aus dem die
Falle selbst besteht.

\begin{figure}
\centering
\captionsetup{width=0.9\textwidth}

\includegraphics[width=0.9\textwidth]{img/abb2.jpg}
\caption{Abbildung 2: Eine Klebefalle zur Überwachung von Papierfischchen und anderen
krabbelnden Insekten.}
\alt{In einer Zimmerecke steht eine Pappfalle in Form eines Trapez auf dem Boden. Sie ist oberseitig beschriftet mit "Monitoring-System" und an der zugewandten Seite mit "without active ingredient".}
\end{figure}

In den Objektsammlungen kommen zusätzlich noch andere Varianten von
Insektenfallen zum Einsatz. Dazu zählen beispielsweise mit Pheromonen
behandelte Klebefallen, die männliche Motten anziehen und damit nicht
nur einen Beitrag zur Überwachung, sondern auch zur Bekämpfung leisten,
da sie den Fortpflanzungszyklus unterbrechen. Für lichtsensible
Fluginsekten werden dagegen Fallen mit einer Lampe ausgebracht, die
blaues UV-Licht aussendet. Die Wahl der Fallenart richtet sich dabei
nach der Art der in den Räumlichkeiten gelagerten Objekte.

Wenn eine Datengrundlage bezüglich der vorhandenen Schädlinge vorliegt,
kann die Diagnose der Einfallswege beginnen. Liegt es beispielsweise an
den geöffneten oder undichten Fenstern? An Besuchenden oder anderen
Personen, die über Kleidung und Gepäck Insekten einschleppen? Wurden
durch Neuzugänge oder deren Verpackungsmaterial Schädlinge in die
Sammlung getragen? Je genauer die Ursachen identifiziert werden, desto
besser können für die Zukunft Vorkehrungen getroffen werden.

Naturgemäß sind eine gute und regelmäßige Reinigung der Räume, ein
angemessenes Raumklima und die Verhinderung der Einführung neuer
Schädlinge (durch Bestandszuwächse oder Ähnliches) die Basis der
Vorsorge. Eingehende Sammlungsobjekte werden bei Ankunft einer genauen
Prüfung unterzogen und auf Schädlinge, aber auch auf Gefahrstoffe hin
untersucht. Idealerweise verbringen sie auch eine gewisse Zeit in
Quarantäne und werden dort gesäubert.

Sollte doch ein Schädlingsbefall festgestellt werden, gibt es heute
viele Bekämpfungsmethoden, die ohne Gift auskommen. Wie erwähnt leistet
bereits das Ausbringen der Klebefallen, deren Anzahl bei Bedarf erhöht
wird, einen gewissen Beitrag zur Bekämpfung. Sind spezifische Objekte
befallen, können sie je nach Größe und Material beispielsweise über
mehrere Tage mit Kälte, Hitze oder Stickstoff behandelt werden. Kleinere
Objekte können dabei in speziellen Schränken oder Räumen untergebracht
werden, während größere in Folie eingewickelt werden. So wurde das oben
erwähnte, von Motten befallene, Auto luftdicht in Alufolie verpackt, die
mit Stickstoff gefüllt wurde. Dies geschah mitten im Publikumsverkehr,
in den Ausstellungsräumen des Museums, um Besuchende über den Vorgang zu
informieren.

Sind Objekte sehr groß -- Flugzeuge, Eisenbahnen und Ähnliches -- können
gegen gewisse Schädlinge auch \enquote{Nützlinge} ausgebracht werden.
Dabei handelt es sich um Insekten (oder ähnliche Kleintiere), die andere
Insekten als Nahrung nutzen. In Fahr- und Flugzeugen des Deutschen
Museums werden so regelmäßig Schlupfwespen (\emph{Trichogramma})
eingesetzt. Diese parasitieren die Eier von Motten, verhindern dadurch
das Entstehen neuer Generationen und sterben in der Folge selbst ab. Ist
der Befall im Gebäude selbst verortet, zum Beispiel im Parkettboden,
besteht auch die Möglichkeit der Bestrahlung mit Mikrowellen.

Im modernen IPM werden also Methoden zur Prävention, zum Monitoring und
zur Bekämpfung auf die vorhandenen Räumlichkeiten, die bewahrten
Sammlungsobjekte und die vorkommenden Schädlinge abgestimmt und
bedarfsgerecht angepasst. Neue Herausforderungen ergeben sich dabei
durch neue Schädlingsarten und Verhaltensänderungen, die durch den
Klimawandel beeinflusst werden können.

\section{Einfluss des Klimawandels auf
Schädlinge}\label{einfluss-des-klimawandels-auf-schuxe4dlinge}

Dass der Klimawandel die Verbreitung von Schädlingen in Museen und
anderen Kultureinrichtungen verändert, ist eine naheliegende Vermutung,
deren konkrete Erforschung sich aber noch in der Anfangsphase
befindet.\footnote{Pascal Querner et al.~(2022). Climate Change and Its
  Effects on Indoor Pests (Insect and Fungi) in Museums. In: Climate
  10(7), 103, \url{https://doi.org/10.3390/cli10070103}} Ausgangspunkt
der in diesem Beitrag aufgegriffenen Sonderausstellung des
Naturhistorischen Museums Wien und des Kurators Pascal Querner war
deshalb auch ein Forschungsprojekt, das genau diese Einflüsse des
Klimawandels in österreichischen Museen empirisch untersucht hat.
Während die Ergebnisse aktuell noch nicht publiziert wurden, sind sie in
Grundzügen schon in der Ausstellung skizziert und bestätigen den
Einfluss des Klimawandels auch auf Schädlingspopulationen in
Kultureinrichtungen.

Das Einschleppen invasiver Arten (Neobiota) in Museen wurde bereits kurz
anhand des Südlichen Nagekäfer thematisiert. Ein weiteres Beispiel ist
der aus Nordamerika eingeschleppte Amerikanische Wespenkäfer
(\emph{Reesa vespulae}), der 2018 in zwei österreichischen Museen
nachgewiesen war, während er sich 2023 bereits in vierzehn Einrichtungen
verbreitet hatte. Während der Klimawandel in der Regel nicht allein für
das Eintreffen dieser Lebewesen verantwortlich ist -- der Mensch und die
Globalisierung leisten hier einen signifikanten Beitrag --, so kann ein
verändertes Klima doch ihre permanente Ansiedlung vor Ort begünstigen.
Denn insbesondere Insekten können ihre Körpertemperatur nicht selbst
steuern und sind deshalb maßgeblich von der Umgebungstemperatur
abhängig. Bei wärmerem Wetter wachsen sie nicht nur schneller heran --
ein Temperaturanstieg von wenigen Grad kann die Entwicklungszeit mehr
als halbieren --, sie zeugen auch mehr Nachkommen, sterben dafür aber
schneller. Wenn die warme Jahreszeit länger wird, können entsprechend im
Jahr mehr Insektengenerationen heranwachsen.

\section{Zum Schluss}\label{zum-schluss}

Insekten und andere Tiere können nicht nur im Haushalt, sondern auch in
Kultureinrichtungen großen Schaden anrichten. Inzwischen gibt es
vielfältige Möglichkeiten, solchen Schäden vorzubeugen und
Schädlingsbefall auf umwelt- und gesundheitsfreundliche Weise zu
bekämpfen. In Museen wird inzwischen aktiv zu diesem Thema mit
Besuchenden kommuniziert, die -- wie die hier besprochene
Sonderausstellung gezeigt hat -- durchaus großes Interesse zeigen. Auch
Bibliotheken können sich überlegen, solche Aspekte in der eigenen
Außenkommunikation zu berücksichtigen, denn insbesondere die durch den
Klimawandel bedingten Veränderungen werden in Zukunft weitere
Anpassungen in der Schädlingsbekämpfung nach sich ziehen und
Auswirkungen auf uns alle haben.

%autor
\begin{center}\rule{0.5\linewidth}{0.5pt}\end{center}

\textbf{Eva Bunge} (\url{https://orcid.org/0000-0002-5587-5934}) ist stellvertretende Leiterin und
Open-Access-Beauftragte des Deutschen Museums in München. Nach dem
Physikstudium wechselte sie ins Bibliothekswesen und arbeitet dort
schwerpunktmäßig zu Open Access und Citizen Science. Sie engagiert sich
bei Library Carpentry und im Beirat der Arbeitsgemeinschaft der
Spezialbibliotheken.

\end{document}
